\chapter{Platform Design}
\label{ch:design}

Where Chapter~\ref{ch:arch} presented the architecture at the level
of boxes and arrows, this chapter describes each module in terms of
its Python types, key responsibilities, and design rationale. Every
reference in this chapter points to a real file in the repository.

\section{Project Structure}
\label{sec:project-structure}

\begin{lstlisting}[style=alBash, caption={Top-level project structure.},
                    label={lst:tree}]
automl-mcp-platform/
|-- app/
|   |-- streamlit_app.py     # Streamlit entrypoint (thin UI)
|   `-- components/          # reserved for render helpers
|-- core/
|   |-- config.py            # pydantic-settings + derived paths
|   |-- logger.py            # named-logger factory
|   `-- workspace.py         # datasets + artifacts registry (local FS)
|-- mcp_servers/
|   |-- data.py              # ingest, list, pandera validation
|   |-- eda.py               # run_full_eda, correlation, prompt
|   |-- modeling.py          # bake-off + Optuna sweep
|   |-- explain.py           # SHAP + native feature importances
|   `-- export.py            # Jupyter notebook + PDF report
|-- orchestrator/
|   |-- pool.py              # MCPClientPool over FastMCPTransport
|   `-- loop.py              # Groq tool-calling loop
|-- data/{uploads,artifacts,tmp}/
|-- requirements.txt
|-- .env.example
`-- README.md
\end{lstlisting}

The choice to keep four sibling top-level packages
(\code{app}, \code{core}, \code{mcp\_servers}, \code{orchestrator})
rather than nesting them under a single \code{src/} package is
deliberate: it makes the dependency graph flat and readable
(\code{app $\rightarrow$ orchestrator $\rightarrow$ mcp\_servers $\rightarrow$ core}),
and it keeps the Streamlit entrypoint invokable with
\code{streamlit run app/streamlit\_app.py} without any wrapper.

\section{Configuration}
\label{sec:config}

Configuration is centralised in \file{core/config.py}:

\begin{lstlisting}[caption={\file{core/config.py} (extract).},
                    label={lst:config}]
class Settings(BaseSettings):
    GROQ_API_KEY: str = ""
    GROQ_MODEL: str = "llama-3.3-70b-versatile"
    MAX_TOOL_ITERATIONS: int = 8
    MAX_UPLOAD_MB: int = 200
    DATA_DIR: str = str(REPO_ROOT / "data")

    @property
    def data_path(self) -> Path: ...
    @property
    def uploads_path(self) -> Path: ...
    @property
    def artifacts_path(self) -> Path: ...
    @property
    def tmp_path(self) -> Path: ...
\end{lstlisting}

Every setting is env-driven through \code{pydantic-settings}
\citep{pydantic}. Derived paths are exposed as properties that
create the directory on first access, so no module has to worry
about directory existence. Appendix~\ref{app:config} lists every
setting with its default and description.

\section{Logging}
\label{sec:logging}

\file{core/logger.py} exposes a single \code{get\_logger(name)}
factory that returns a name-scoped logger with a compact stderr
formatter. This keeps log output consistent across every module and
avoids the common ``duplicate handler'' problem when Streamlit
re-runs the script:

\begin{lstlisting}[caption={\file{core/logger.py}.},label={lst:logger}]
def get_logger(name: str) -> logging.Logger:
    logger = logging.getLogger(name)
    if logger.handlers:
        return logger  # idempotent
    logger.setLevel(logging.INFO)
    handler = logging.StreamHandler(sys.stderr)
    handler.setFormatter(logging.Formatter(
        fmt="%(asctime)s | %(levelname)-7s | %(name)s | %(message)s",
        datefmt="%H:%M:%S"))
    logger.addHandler(handler)
    logger.propagate = False
    return logger
\end{lstlisting}

\section{The Workspace Module}
\label{sec:workspace}

\file{core/workspace.py} is the sole persistence layer. It contains
two dataclasses --- \code{DatasetRecord} and \code{ArtifactRecord}
--- and a small set of pure functions.

\paragraph{Dataset operations.}
\code{save\_uploaded\_dataset(filename, data, \dots)} writes raw
bytes into \file{data/uploads/}, deduplicating filenames by appending
a Unix timestamp when necessary, and returns a \code{DatasetRecord}
with a UUID and populated row/column metadata (if supplied by the
caller). \code{list\_datasets()} enumerates the uploads directory and
returns dictionaries sorted by modification time.
\code{resolve\_dataset(file\_ref)} accepts either an absolute path or
a base filename and returns the resolved \code{Path}, raising
\code{FileNotFoundError} with a helpful message when the file is not
present.

\paragraph{Artefact operations.}
\code{save\_artifact(local\_path, kind, metadata=\dots, move=True)}
is the single entry point that every MCP tool uses. It:
\begin{enumerate}[leftmargin=1.4em]
  \item Resolves the target category from a kind-to-category
        dictionary (\code{plot} $\rightarrow$ \code{plots},
        \code{model} $\rightarrow$ \code{models}, etc.).
  \item Assigns a UUID and constructs a destination path of the form
        \code{data/artifacts/<category>/<uuid[:8]>\_\_<safe\_name>}.
  \item Moves or copies the file, computes MIME type via
        \code{mimetypes}, and appends a JSON record to
        \file{data/artifacts/\_index.json}.
\end{enumerate}
The registry is a plain JSON array to keep the code trivial;
switching to SQLite or an experiment tracker (Section~\ref{sec:future-mlops})
requires changing only this file.

\section{The MCP Servers}
\label{sec:servers-design}

Each server module follows the same template:

\begin{lstlisting}[caption={MCP server module template.},
                    label={lst:server-template}]
mcp = FastMCP(name="mcp-<service>")

@mcp.tool
def <tool_name>(arg1: <T1>, arg2: <T2> = <default>, ...) -> dict[str, Any]:
    """Short description used as the LLM-visible tool description.

    Args:
        arg1: ...
        arg2: ...
    """
    ...
    return {"...": ...}
\end{lstlisting}

Tools whose work is compute-bound or IO-bound are declared
\code{async def} and additionally receive a \code{ctx: Context}
parameter through which they emit progress with
\code{ctx.report\_progress(current, total, message)}.

\subsection{\svc{mcp-data}}
Three tools: \tool{list\_uploaded\_files},
\tool{ingest\_dataset}, \tool{validate\_schema\_with\_pandera}.
The validation tool infers a \code{pandera.DataFrameSchema} by
walking each column: numeric columns receive an \code{in\_range} check,
low-cardinality string columns receive an \code{isin} check, and
nullability follows from observed NaN values. The tool returns the
first 50 failure cases (column, check, offending value, index) to
make debugging tractable.

\subsection{\svc{mcp-eda}}
Two tools: \tool{run\_full\_eda} and \tool{render\_correlation\_matrix},
and one prompt: \tool{eda\_deep\_dive}. The full EDA tool renders
three PNGs into a per-service scratch directory
(\file{data/tmp/eda/}) before handing them to
\code{save\_artifact}. Correlation entries above $|\rho|\ge 0.6$ are
returned as ``strong correlations'' so that the LLM can surface them
in prose.

\subsection{\svc{mcp-modeling}}
Two tools: \tool{run\_parallel\_bake\_off} and
\tool{trigger\_hyperparameter\_sweep}.
The bake-off constructs a shared
\code{ColumnTransformer} preprocessor and spawns a
\code{ThreadPoolExecutor} with up to four workers, one per
candidate. Each worker runs $k$-fold cross-validation with either
\code{f1\_weighted} or \code{r2}, then refits on the full training
split and evaluates on the held-out test split. Champion selection
uses the primary metric (F1 for classification, $R^2$ for
regression). The champion pipeline is serialised with \code{joblib}
along with the feature columns, problem type, and target column, and
persisted as a \code{model} artefact. A 500-row sample of
\code{X\_train} is written as parquet and persisted as an
\code{xtrain\_sample} artefact for downstream SHAP.

\subsection{\svc{mcp-explain}}
Two tools: \tool{calculate\_shap\_values} and
\tool{generate\_feature\_importance\_plot}.
Both tools read the champion artefact by ID via
\code{workspace.artifact\_path}, so the ID plumbed from the bake-off
result is the only coupling between the modelling and explainability
stages.

\subsection{\svc{mcp-export}}
Two tools: \tool{generate\_jupyter\_notebook} and
\tool{compile\_pdf\_report}. Notebooks are built as pure Python
dictionaries following the nbformat schema and serialised to disk
with \code{json.dumps}. PDFs are built with ReportLab
\citep{reportlab}: a title block, a summary paragraph, a leaderboard
table, and an optional notes section.

\section{Orchestration Design}
\label{sec:orch-design}

Two responsibilities live in \file{orchestrator/}: aggregating the
tools into a single catalog (\file{pool.py}) and iterating on the
LLM's tool-call loop (\file{loop.py}).

\subsection{The Pool}
\label{sec:pool-design}

\code{MCPClientPool} maintains three registries: \code{tool\_index}
(tool name $\to$ owning service name), \code{tool\_schemas}
(OpenAI-compatible schema list ready to be passed to Groq), and
\code{prompt\_index}. Registries are rebuilt on start-up and can be
refreshed after a service reconnect.

A key design decision is name-collision handling. If two servers
register a tool with the same name, the second one is registered
under a canonical name of the form
\code{<service>\_\_<original\_name>}. This is deliberately loud so
that the developer notices and renames.

\subsection{The Loop}
\label{sec:loop-design}

\code{stream\_chat} is an asynchronous generator of event dicts.
Each iteration of the outer \code{for} loop is one Groq streaming
call. The loop accumulates \code{tool\_calls} chunks, materialises
them into a list, executes them (with per-tool progress bridged
through an \code{asyncio.Queue}), appends the tool results back to
the message history, and re-invokes the LLM. The loop terminates
either when Groq emits no \code{tool\_calls} in an iteration or when
\code{MAX\_TOOL\_ITERATIONS} is reached.

The generator emits six event kinds --- \code{token},
\code{tool\_start}, \code{tool\_progress}, \code{tool\_end},
\code{done}, \code{error} --- documented in
Section~\ref{sec:event-schema}. This design is intentionally
UI-agnostic: any consumer (CLI, Jupyter widget, Streamlit,
FastAPI-over-SSE) can subscribe.

\section{Streamlit Front-end}
\label{sec:st-design}

The UI has three tabs: \emph{Chat}, \emph{Artifacts}, and
\emph{Tools}. A sidebar contains the upload widget, the active
dataset selector, and a live service-status list built from
\code{pool.snapshot()}.

To bridge Streamlit's synchronous execution model with the async
orchestrator, the UI creates a single \code{PoolRunner} object cached
via \code{st.cache\_resource}. \code{PoolRunner} owns a dedicated
event loop running in a daemon thread and exposes a
\code{stream\_events(query, active\_file, history)} method that
returns a \code{queue.Queue}. Each user turn consumes the queue in
the main thread and updates Streamlit placeholders in place. This
gives the appearance of real-time streaming without pushing async
into the Streamlit script.

\section{Extensibility}
\label{sec:extensibility}

Adding a new capability follows a strict recipe:

\begin{enumerate}[leftmargin=1.4em]
  \item Add a \code{@mcp.tool}-decorated function to the appropriate
        \file{mcp\_servers/<service>.py} (or create a new server
        module and register it in
        \file{orchestrator/pool.py::register\_default\_servers}).
  \item If the tool generates a new kind of file, add a
        kind-to-category entry to \code{\_KIND\_TO\_CATEGORY} in
        \file{core/workspace.py}.
  \item Optionally, add a friendly label to \code{TOOL\_LABELS} in
        \file{orchestrator/loop.py}.
\end{enumerate}

No change is required to the Streamlit UI or to the orchestration
loop.
